% GENERATED FILE. Edit canonical lessons, then run npm run generate:books.
% canonical-source-hash: fnv1a64:d883cd20d9aee91b
% canonical-lessons: SA-C392-abhyarthayate, SA-C392-sahayakah, SA-C392-asvasayati, SA-C392-sahakarah, SA-C392-anugrahah

\chapter{Asking for Help}
\label{ch:sa-asking-for-help}
\hlchaptermodality{392}

\begin{chapteropening}
\textbf{By the end of this chapter:} \emph{I can say who will help if I ask.}

Complete the last lesson of chapter 392: I can say who will help if I ask.
\end{chapteropening}

\section[\texorpdfstring{\sk{अभ्यर्थयते} (abhyarthayate)}{अभ्यर्थयते (abhyarthayate)}]{\sk{अभ्यर्थयते} (abhyarthayate) — requests; asks for}

\label{lesson:SA-C392-abhyarthayate}



\begin{quote}
Before the new one: say the Sanskrit for suitable, then the Sanskrit for spoils.
\end{quote}

\subsection*{You'll want to know: \sk{अभ्यर्थयते}}
\textbf{\sk{अभ्यर्थयते}} — \emph{abhyarthayate} — \textquotedblleft{}requests; asks for\textquotedblright{}.

\begin{quote}
\sk{यदि} \sk{भवान्} \sk{अभ्यर्थयते}, \sk{तर्हि} \sk{सः} \sk{साहाय्यं} \sk{करिष्यति।}
\end{quote}

\begin{quote}
\emph{yadi bhavān abhyarthayate, tarhi saḥ sāhāyyaṃ kariṣyati.}
\end{quote}

\begin{quote}
If you ask him, he will help.
\end{quote}

Past (he …): \sk{सः} \sk{साहाय्यम्} \sk{अभ्यर्थितवान्}, \emph{saḥ sāhāyyam abhyarthitavān}.

\begin{cousinweb}
\textbf{\sk{अभ्यर्थयते}} (\emph{abhyarthayate}) is built on \emph{abhi}, \textquotedblleft{}towards\textquotedblright{}, and the verb \emph{arthaya}, \textquotedblleft{}to ask, to wish\textquotedblright{}; MW gives \emph{abhy-√arth} \textquotedblleft{}to request, to ask for\textquotedblright{} and \emph{abhy-arthita} \textquotedblleft{}asked, invited\textquotedblright{} (MW).
\end{cousinweb}

\begin{grammarlens}[title={What you've built}]
One more word for saying what happens if.
\end{grammarlens}

\subsection*{Your turn}
\begin{itemize}
\raggedright
  \item \emph{Say it:} \emph{abhyarthayate}
  \item \emph{Say it:} \emph{abhyarthayate}, once more
  \item \emph{Say it:} \emph{yadi bhavān abhyarthayate, tarhi saḥ sāhāyyaṃ kariṣyati.} — If you ask him, he will help.
  \item \emph{Recall:} say the Sanskrit for suitable, then the Sanskrit for spoils, then say \emph{abhyarthayate} again
\end{itemize}

\begin{tcolorbox}[breakable,colback=teal!4,colframe=teal!35!black,title={Before you move on}]
What does \sk{अभ्यर्थयते} mean? (\textquotedblleft{}Requests; asks for\textquotedblright{}.) Say it once more.
\end{tcolorbox}
\section[\texorpdfstring{\sk{सहायकः} (sahāyakaḥ)}{सहायकः (sahāyakaḥ)}]{\sk{सहायकः} (sahāyakaḥ) — an assistant; a helper}

\label{lesson:SA-C392-sahayakah}



\begin{quote}
Before the new one: say the Sanskrit for spoils, then the Sanskrit for requests.
\end{quote}

\subsection*{You'll want to know: \sk{सहायकः}}
\textbf{\sk{सहायकः}} — \emph{sahāyakaḥ} — \textquotedblleft{}an assistant; a helper\textquotedblright{}.

\begin{quote}
\sk{यदि} \sk{कार्यं} \sk{बहु}, \sk{तर्हि} \sk{एकं} \sk{सहायकम्} \sk{आनयतु।}
\end{quote}

\begin{quote}
\emph{yadi kāryaṃ bahu, tarhi ekaṃ sahāyakam ānayatu.}
\end{quote}

\begin{quote}
If there is a lot of work, bring an assistant.
\end{quote}

A masculine noun.

\begin{cousinweb}
\textbf{\sk{सहायकः}} (\emph{sahāyakaḥ}) is built on \emph{sahāya}, \textquotedblleft{}a companion\textquotedblright{}; MW gives \emph{sahāyaka} as the same as \emph{sahāya}, a companion who helps (MW).
\end{cousinweb}

\begin{grammarlens}[title={What you've built}]
One more word for saying what happens if.
\end{grammarlens}

\subsection*{Your turn}
\begin{itemize}
\raggedright
  \item \emph{Say it:} \emph{sahāyakaḥ}
  \item \emph{Say it:} \emph{sahāyakaḥ}, once more
  \item \emph{Say it:} \emph{yadi kāryaṃ bahu, tarhi ekaṃ sahāyakam ānayatu.} — If there is a lot of work, bring an assistant.
  \item \emph{Recall:} say the Sanskrit for spoils, then the Sanskrit for requests, then say \emph{sahāyakaḥ} again
\end{itemize}

\begin{tcolorbox}[breakable,colback=teal!4,colframe=teal!35!black,title={Before you move on}]
What does \sk{सहायकः} mean? (\textquotedblleft{}An assistant; a helper\textquotedblright{}.) Say it once more.
\end{tcolorbox}
\section[\texorpdfstring{\sk{आश्वासयति} (āśvāsayati)}{आश्वासयति (āśvāsayati)}]{\sk{आश्वासयति} (āśvāsayati) — reassures; comforts}

\label{lesson:SA-C392-asvasayati}



\begin{quote}
Before the new one: say the Sanskrit for requests, then the Sanskrit for an assistant.
\end{quote}

\subsection*{You'll want to know: \sk{आश्वासयति}}
\textbf{\sk{आश्वासयति}} — \emph{āśvāsayati} — \textquotedblleft{}reassures; comforts\textquotedblright{}.

\begin{quote}
\sk{यदि} \sk{बालः} \sk{रोदिति}, \sk{तर्हि} \sk{माता} \sk{तम्} \sk{आश्वासयति।}
\end{quote}

\begin{quote}
\emph{yadi bālaḥ roditi, tarhi mātā tam āśvāsayati.}
\end{quote}

\begin{quote}
If the child cries, its mother comforts it.
\end{quote}

Past (she …): \sk{सा} \sk{बालम्} \sk{आश्वासितवती}, \emph{sā bālam āśvāsitavatī}.

\begin{cousinweb}
\textbf{\sk{आश्वासयति}} (\emph{āśvāsayati}) is built on \emph{ā} and the root √śvas, \textquotedblleft{}to breathe\textquotedblright{}, in its causative form; MW gives \emph{ā-śvāsita} \textquotedblleft{}encouraged, comforted, consoled\textquotedblright{} (MW).
\end{cousinweb}

\begin{grammarlens}[title={What you've built}]
One more word for saying what happens if.
\end{grammarlens}

\subsection*{Your turn}
\begin{itemize}
\raggedright
  \item \emph{Say it:} \emph{āśvāsayati}
  \item \emph{Say it:} \emph{āśvāsayati}, once more
  \item \emph{Say it:} \emph{yadi bālaḥ roditi, tarhi mātā tam āśvāsayati.} — If the child cries, its mother comforts it.
  \item \emph{Recall:} say the Sanskrit for requests, then the Sanskrit for an assistant, then say \emph{āśvāsayati} again
\end{itemize}

\begin{tcolorbox}[breakable,colback=teal!4,colframe=teal!35!black,title={Before you move on}]
What does \sk{आश्वासयति} mean? (\textquotedblleft{}Reassures; comforts\textquotedblright{}.) Say it once more.
\end{tcolorbox}
\section[\texorpdfstring{\sk{सहकारः} (sahakāraḥ)}{सहकारः (sahakāraḥ)}]{\sk{सहकारः} (sahakāraḥ) — co-operation}

\label{lesson:SA-C392-sahakarah}



\begin{quote}
Before the new one: say the Sanskrit for an assistant, then the Sanskrit for reassures.
\end{quote}

\subsection*{You'll want to know: \sk{सहकारः}}
\textbf{\sk{सहकारः}} — \emph{sahakāraḥ} — \textquotedblleft{}co-operation\textquotedblright{}.

\begin{quote}
\sk{यदि} \sk{सहकारः} \sk{अस्ति}, \sk{तर्हि} \sk{कठिनं} \sk{कार्यम्} \sk{अपि} \sk{सुकरम्।}
\end{quote}

\begin{quote}
\emph{yadi sahakāraḥ asti, tarhi kaṭhinaṃ kāryam api sukaram.}
\end{quote}

\begin{quote}
If there is co-operation, even hard work is easy.
\end{quote}

A masculine noun.

\begin{cousinweb}
\textbf{\sk{सहकारः}} (\emph{sahakāraḥ}) is built from \emph{saha}, \textquotedblleft{}with\textquotedblright{}, and \emph{kāra}, \textquotedblleft{}a doing\textquotedblright{}; MW gives \emph{saha-kāra} \textquotedblleft{}acting with, co-operation, assistance\textquotedblright{} (MW).
\end{cousinweb}

\begin{grammarlens}[title={What you've built}]
One more word for saying what happens if.
\end{grammarlens}

\subsection*{Your turn}
\begin{itemize}
\raggedright
  \item \emph{Say it:} \emph{sahakāraḥ}
  \item \emph{Say it:} \emph{sahakāraḥ}, once more
  \item \emph{Say it:} \emph{yadi sahakāraḥ asti, tarhi kaṭhinaṃ kāryam api sukaram.} — If there is co-operation, even hard work is easy.
  \item \emph{Recall:} say the Sanskrit for an assistant, then the Sanskrit for reassures, then say \emph{sahakāraḥ} again
\end{itemize}

\begin{tcolorbox}[breakable,colback=teal!4,colframe=teal!35!black,title={Before you move on}]
What does \sk{सहकारः} mean? (\textquotedblleft{}Co-operation\textquotedblright{}.) Say it once more.
\end{tcolorbox}
\section[\texorpdfstring{\sk{अनुग्रहः} (anugrahaḥ)}{अनुग्रहः (anugrahaḥ)}]{\sk{अनुग्रहः} (anugrahaḥ) — a kindness; a favour}

\label{lesson:SA-C392-anugrahah}



\begin{quote}
Before the new one: say the Sanskrit for reassures, then the Sanskrit for co-operation.
\end{quote}

\subsection*{You'll want to know: \sk{अनुग्रहः}}
\textbf{\sk{अनुग्रहः}} — \emph{anugrahaḥ} — \textquotedblleft{}a kindness; a favour\textquotedblright{}.

\begin{quote}
\sk{यदि} \sk{अध्यापकस्य} \sk{अनुग्रहः} \sk{अस्ति}, \sk{तर्हि} \sk{ज्ञानं} \sk{सुलभं} \sk{भवति।}
\end{quote}

\begin{quote}
\emph{yadi adhyāpakasya anugrahaḥ asti, tarhi jñānaṃ sulabhaṃ bhavati.}
\end{quote}

\begin{quote}
If the teacher is kind, knowledge comes easily.
\end{quote}

A masculine noun.

\begin{cousinweb}
\textbf{\sk{अनुग्रहः}} (\emph{anugrahaḥ}) is built on \emph{anu}, \textquotedblleft{}after\textquotedblright{}, and the root √grah, \textquotedblleft{}to seize, to take\textquotedblright{}; MW gives \emph{anu-graha} \textquotedblleft{}favour, kindness, conferring benefits\textquotedblright{} (MW).
\end{cousinweb}

\begin{grammarlens}[title={What you've built}]
One more word for saying what happens if.
\end{grammarlens}

\subsection*{Your turn}
\begin{itemize}
\raggedright
  \item \emph{Say it:} \emph{anugrahaḥ}
  \item \emph{Say it:} \emph{anugrahaḥ}, once more
  \item \emph{Say it:} \emph{yadi adhyāpakasya anugrahaḥ asti, tarhi jñānaṃ sulabhaṃ bhavati.} — If the teacher is kind, knowledge comes easily.
  \item \emph{Recall:} say the Sanskrit for reassures, then the Sanskrit for co-operation, then say \emph{anugrahaḥ} again
\end{itemize}

\begin{tcolorbox}[breakable,colback=teal!4,colframe=teal!35!black,title={Before you move on}]
What does \sk{अनुग्रहः} mean? (\textquotedblleft{}A kindness; a favour\textquotedblright{}.) Say it once more.
\end{tcolorbox}
